\section{Functoriality for quasi-coherent modules}
\label{section-quasi-coherent}

\noindent
Let $X$ be a scheme. We denote $\QCoh(\mathcal{O}_X)$ the
category of quasi-coherent $\mathcal{O}_X$-modules as defined in
Modules, Definition \ref{modules-definition-quasi-coherent}.
We have seen in
Section \ref{section-quasi-coherent-affine}
that the category $\QCoh(\mathcal{O}_X)$ has a lot of good properties
when $X$ is affine. Since the property of being quasi-coherent is
local on $X$, these properties are inherited by the category
of quasi-coherent sheaves on any scheme $X$. We enumerate them here.
\begin{enumerate}
\item A sheaf of $\mathcal{O}_X$-modules $\mathcal{F}$ is
quasi-coherent if and only if the restriction of $\mathcal{F}$
to each affine open $U = \Spec(R)$ is of the form
$\widetilde M$ for some $R$-module $M$.
\item A sheaf of $\mathcal{O}_X$-modules $\mathcal{F}$ is
quasi-coherent if and only if the restriction of $\mathcal{F}$
to each of the members of an affine open covering is quasi-coherent.
\item Any direct sum of quasi-coherent sheaves is quasi-coherent.
\item Any colimit of quasi-coherent sheaves is quasi-coherent.
\item
\label{item-kernel}
The kernel and cokernel of a morphism of quasi-coherent sheaves
is quasi-coherent.
\item Given a short exact sequence of $\mathcal{O}_X$-modules
$0 \to \mathcal{F}_1 \to \mathcal{F}_2 \to \mathcal{F}_3 \to 0$
if two out of three are quasi-coherent so is the third.
\item Given a morphism of schemes $f : Y \to X$ the pullback
of a quasi-coherent $\mathcal{O}_X$-module is a quasi-coherent
$\mathcal{O}_Y$-module. See
Modules, Lemma \ref{modules-lemma-pullback-quasi-coherent}.
\item Given two quasi-coherent $\mathcal{O}_X$-modules
the tensor product is quasi-coherent, see
Modules, Lemma \ref{modules-lemma-tensor-product-permanence}.
\item Given a quasi-coherent $\mathcal{O}_X$-module $\mathcal{F}$
the tensor, symmetric and exterior algebras on $\mathcal{F}$
are quasi-coherent, see
Modules, Lemma \ref{modules-lemma-whole-tensor-algebra-permanence}.
\item Given two quasi-coherent $\mathcal{O}_X$-modules
$\mathcal{F}$, $\mathcal{G}$ such that $\mathcal{F}$
is of finite presentation, then the internal hom
$\SheafHom_{\mathcal{O}_X}(\mathcal{F}, \mathcal{G})$
is quasi-coherent, see
Modules, Lemma \ref{modules-lemma-internal-hom-locally-kernel-direct-sum}
and (\ref{item-kernel}) above.
\end{enumerate}
On the other hand, it is in general not the case that the
pushforward of a quasi-coherent module is quasi-coherent.
Here is a case where this does hold.

\begin{lemma}
\label{lemma-push-forward-quasi-coherent}
Let $f : X \to S$ be a morphism of schemes.
If $f$ is quasi-compact and quasi-separated then
$f_*$ transforms quasi-coherent $\mathcal{O}_X$-modules
into quasi-coherent $\mathcal{O}_S$-modules.
\end{lemma}

\begin{proof}
The question is local on $S$ and hence we may assume that
$S$ is affine. Because $X$ is quasi-compact we may write
$X = \bigcup_{i = 1}^n U_i$ with each $U_i$ open affine.
Because $f$ is quasi-separated we may write
$U_i \cap U_j = \bigcup_{k = 1}^{n_{ij}} U_{ijk}$ for some
affine open $U_{ijk}$, see Lemma \ref{lemma-characterize-quasi-separated}.
Denote $f_i : U_i \to S$ and $f_{ijk} : U_{ijk} \to S$ the
restrictions of $f$. For any open $V$ of $S$ and any sheaf
$\mathcal{F}$ on $X$ we have
\begin{eqnarray*}
f_*\mathcal{F}(V) & = & \mathcal{F}(f^{-1}V) \\
& = &
\Ker\left(
\bigoplus\nolimits_i \mathcal{F}(f^{-1}V \cap U_i)
\to
\bigoplus\nolimits_{i, j, k} \mathcal{F}(f^{-1}V \cap U_{ijk})\right) \\
& = &
\Ker\left(
\bigoplus\nolimits_i f_{i, *}(\mathcal{F}|_{U_i})(V)
\to
\bigoplus\nolimits_{i, j, k} f_{ijk, *}(\mathcal{F}|_{U_{ijk}})(V)\right) \\
& = &
\Ker\left(
\bigoplus\nolimits_i f_{i, *}(\mathcal{F}|_{U_i})
\to
\bigoplus\nolimits_{i, j, k} f_{ijk, *}(\mathcal{F}|_{U_{ijk}})\right)(V)
\end{eqnarray*}
In other words there is an exact sequence of sheaves
$$
0 \to f_*\mathcal{F}
\to \bigoplus f_{i, *}\mathcal{F}_i
\to \bigoplus f_{ijk, *}\mathcal{F}_{ijk}
$$
where $\mathcal{F}_i, \mathcal{F}_{ijk}$ denotes the
restriction of $\mathcal{F}$ to the corresponding open.
If $\mathcal{F}$ is a quasi-coherent $\mathcal{O}_X$-module
then $\mathcal{F}_i$ is a quasi-coherent $\mathcal{O}_{U_i}$-module
and $\mathcal{F}_{ijk}$ is a quasi-coherent $\mathcal{O}_{U_{ijk}}$-module.
Hence by Lemma \ref{lemma-widetilde-pullback} we see that the second and
third term of the exact sequence are quasi-coherent
$\mathcal{O}_S$-modules. Thus we conclude that
$f_*\mathcal{F}$ is a quasi-coherent $\mathcal{O}_S$-module.
\end{proof}

\noindent
Using this we can characterize (closed) immersions of schemes
as follows.

\begin{lemma}
\label{lemma-characterize-closed-immersions}
Let $f : X \to Y$ be a morphism of schemes.
Suppose that
\begin{enumerate}
\item $f$ induces a homeomorphism of $X$ with a
closed subset of $Y$, and
\item $f^\sharp : \mathcal{O}_Y \to f_*\mathcal{O}_X$
is surjective.
\end{enumerate}
Then $f$ is a closed immersion of schemes.
\end{lemma}

\begin{proof}
Assume (1) and (2). By (1) the morphism $f$ is quasi-compact
(see Topology, Lemma \ref{topology-lemma-closed-in-quasi-compact}).
Conditions (1) and (2) imply conditions (1) and (2) of
Lemma \ref{lemma-injective-points-surjective-stalks}.
Hence $f : X \to Y$ is a monomorphism. In particular, $f$
is separated, see Lemma \ref{lemma-monomorphism-separated}.
Hence Lemma \ref{lemma-push-forward-quasi-coherent}
above applies and we conclude that $f_*\mathcal{O}_X$ is a quasi-coherent
$\mathcal{O}_Y$-module. Therefore the kernel of
$\mathcal{O}_Y \to f_*\mathcal{O}_X$ is quasi-coherent
by Lemma \ref{lemma-extension-quasi-coherent}. Since a quasi-coherent
sheaf is locally generated by sections (see
Modules, Definition \ref{modules-definition-quasi-coherent})
this implies that $f$ is a closed immersion, see
Definition \ref{definition-closed-immersion-locally-ringed-spaces}.
\end{proof}

\noindent
We can use this lemma to prove the following lemma.

\begin{lemma}
\label{lemma-composition-immersion}
A composition of immersions of schemes is an immersion,
a composition of closed immersions of schemes is a closed immersion, and
a composition of open immersions of schemes is an open immersion.
\end{lemma}

\begin{proof}
This is clear for the case of open immersions since an open subspace of
an open subspace is also an open subspace.

\medskip\noindent
Suppose $a : Z \to Y$ and $b : Y \to X$ are closed immersions of schemes.
We will verify that $c = b \circ a$ is also a closed immersion.
The assumption implies that $a$ and $b$ are homeomorphisms onto closed subsets,
and hence also $c = b \circ a$ is a homeomorphism onto a closed subset.
Moreover, the map $\mathcal{O}_X \to c_*\mathcal{O}_Z$ is surjective
since it factors as the composition of the surjective maps
$\mathcal{O}_X \to b_*\mathcal{O}_Y$ and
$b_*\mathcal{O}_Y \to b_*a_*\mathcal{O}_Z$ (surjective as $b_*$ is exact,
see Modules, Lemma \ref{modules-lemma-i-star-exact}).
Hence by Lemma \ref{lemma-characterize-closed-immersions}
above $c$ is a closed immersion.

\medskip\noindent
Finally, we come to the case of immersions.
Suppose $a : Z \to Y$ and $b : Y \to X$ are immersions of schemes.
This means there exist open subschemes
$V \subset Y$ and $U \subset X$ such that
$a(Z) \subset V$, $b(Y) \subset U$ and $a : Z \to V$
and $b : Y \to U$ are closed immersions.
Since the topology on $Y$ is induced from the topology on
$U$ we can find an open $U' \subset U$ such that
$V = b^{-1}(U')$. Then we see that
$Z \to V = b^{-1}(U') \to U'$ is a composition of
closed immersions and hence a closed immersion.
This proves that $Z \to X$ is an immersion and we win.
\end{proof}











